\par\medskip\noindent\textbf{Cómo leer la salida.}
{\small\setlength{\tabcolsep}{4pt}\renewcommand{\arraystretch}{1.15}
\begin{xltabular}{\linewidth}{@{}>{\raggedright\arraybackslash}p{0.24\linewidth}>{\raggedleft\arraybackslash}p{0.15\linewidth}>{\raggedright\arraybackslash}X@{}}
\toprule Indicador & Valor & Interpretación \\ \midrule \endfirsthead
\toprule Indicador & Valor & Interpretación \\ \midrule \endhead
\bottomrule \endlastfoot
Chow (quiebre 1992-10) & $F=11{,}43$; $p=1{,}3\times10^{-5}$ & Rechaza que alfa y beta sean iguales antes y después de octubre de 1992 (bajo homocedasticidad). \\
Beta antes / después & 0,629 / 0,384 & Tras la desregulación, el sector se volvió aún menos sensible al mercado. \\
Alfa anual antes / después & 2,19 / 4,35\,\% & La alfa casi se duplica en el segundo tramo. \\
Wald robusto (HC1) & $W=10{,}97$; $p=0{,}004$ & La versión robusta a heterocedasticidad también rechaza: el quiebre no es un artefacto de la varianza cambiante. \\
sup-$F$ & 16,69 en 1987-02 & Mayor estadístico $F$ al recorrer todas las fechas posibles; supera el valor crítico de Andrews (5,85 en escala $F$), así que hay un quiebre, aunque la fecha más probable es 1987 y no 1992. \\
CUSUM & 0,73 de la banda & La suma acumulada de residuos recursivos nunca cruza las bandas: la prueba CUSUM, de menor potencia ante cambios de pendiente, no detecta inestabilidad. \\
\end{xltabular}}
